\section{Conclusion}

The project has progressed from a finite coordinatewise checker to a family of interfaces that retain the observations required by a consumer and compose them across an entire word. Forced rows provide a concrete connection to the first paper: supplied atomic values determine the tables actually used by the proof transitions, while kernels distinguish local information loss from collapse of the protected carrier. Finite gluing connects those actions to arbitrary lengths once the numerical and source bridges are supplied.

The signed-bound case study exhibits both strength and limits. A shared-prefix carrier supports an exact upper maximum; a two-class carry factor supports a conditional lower-preservation argument. The lower result need not be exact, and its sign condition matters. The Lean pilot now checks the ascending mathematical fragment and the modeled cell-factor bridge.

The main open step is to make adequate observation selection and joint-carrier composition more general, then connect them to a stable implementation and stronger source verification. The present results provide concrete mechanisms, replayable evidence, countermodels to inadequate interfaces, and a machine-checked starting point for that work.
